\section{Conclusion}
\label{sec:conclusion}

We measured three retrieval architectures over one corpus and one question set:
167{,}050 SEC filings from 10{,}511 companies, 508{,}714 embedded chunks and a
graph of 430{,}257 nodes and 561{,}144 edges, queried with 20 questions in four
shape categories, each asked three times, for 180 attempts scored
programmatically rather than by an LLM. GraphRAG solved 14 of 20 questions on all
three attempts, text-to-SQL 9 and vector RAG 3.

The aggregate is the least useful number in the paper. \textbf{GraphRAG lost two
of the four categories}: structured aggregation to text-to-SQL as predicted, at 24
times the token cost, and cross-document entity resolution to text-to-SQL against
prediction. The defensible summary is therefore not that graphs beat vectors but a
rule about question shape:

\begin{itemize}
  \item If the answer lies in one document and the entity is named, the binding
        constraint is \emph{entity filtering}, not ranking. Similarity search
        alone fails at scale --- 3 of 15 here --- and needs either a metadata
        filter on a stable identifier or a traversal to supply one.
  \item If the question is a \texttt{GROUP BY} over typed columns, use SQL. A
        traversal adds no information and adds a place to go wrong.
  \item If the question requires resolving \emph{who} an entity is, or joining
        facts that separate documents assert, structural retrieval is worth its
        cost. This is where the margin was unambiguous: 15 of 15 against 0 of 15
        for both baselines.
  \item If provenance must be auditable, the graph's advantage is categorical
        rather than marginal. It was never uncited across 60 attempts, because the
        traversal returns accession numbers before any text is read, whereas
        text-to-SQL was right-but-unciteable 16 times.
\end{itemize}

The trustworthiness finding is narrower but sharper than the accuracy finding.
Only 2 of 180 attempts asserted something verifiably false, and both were the same
error: text-to-SQL conflating distinct people who share a surname, matching on the
name the question supplied because that is what the text contains. Resolving
identity to an identifier before reading any text prevents that class of error
structurally rather than probabilistically. For an application where a confident
wrong answer costs more than a refusal, that property may matter more than the
pass rate.

We claim none of this beyond one corpus, one model and twenty questions.

\subsection*{Future work}

Four steps would materially strengthen these results. \emph{Re-run all 180
attempts per model}, turning the cross-model probe of \Cref{tab:crossmodel} into a
measurement and separating architectural effects from model effects. \emph{Repair
the entity-resolution layer} --- 0.12\% \texttt{:RESOLVES\_TO} coverage and the
name-normalization mismatch --- and re-measure T4, which is the only way to settle
whether that loss is architectural or incidental. \emph{Obtain an independent
question set}, since the present one was written by the author of the graph, which
is the study's most serious threat. And \emph{vary the retriever}: a longer-context
encoder with a metadata filter on CIK would test directly whether the T1 collapse
is a property of embeddings or of this particular configuration --- our own
diagnosis predicts the metadata filter, not the encoder, does the work.
